\chapter{Introduction}
\label{ch:introduction}
\newpage

\section{Research Context}
\label{sec:research-context}

AI tools for software development have changed in what they do. GitHub Copilot (2021) and Amazon
CodeWhisperer (2022) completed the next line or block, and the developer decided what to keep.
Newer tools work at the level of a whole task. Claude Code (Anthropic, 2025), Cursor's agent mode,
Devin (Cognition, 2024), OpenAI Codex CLI (2025), Google Jules (2025) and open-source tools such as
OpenCode and Aider read an issue, plan steps, edit files, run tests and revise their changes without
a human approving each step. This loop, in which a language model alternates reasoning with tool
calls, was described as the ReAct pattern \citep{yao2023react}. As Figure~\ref{fig:timeline} shows,
assistants (above the axis) suggest code in the editor, while agents (below the axis) take a task and
return a finished change. The first half of 2025 is enlarged because the three agents studied in this
thesis (in gold) were released in that period.

\begin{figure}[htb]
\centering
% Figure: Evolution of AI coding tools (timeline)
\begin{tikzpicture}[
  every node/.style={font=\footnotesize},
  yr/.style={draw=issatblue,thick,rounded corners=2pt,fill=lightblue,minimum width=2.55cm,
             minimum height=0.95cm,align=center,text=issatblue,font=\footnotesize\bfseries},
  ya/.style={draw=issatgold,thick,rounded corners=2pt,fill=issatgold!18,minimum width=2.7cm,
             minimum height=0.95cm,align=center,text=issatblue!90!black,font=\footnotesize\bfseries},
]
\draw[-{Latex[length=2.5mm]},line width=1pt,accentgray] (-0.3,0) -- (13.9,0);
\node[anchor=west,font=\scriptsize\itshape,accentgray] at (13.95,0) {time};

\foreach \x/\yr in {0.6/2021, 3.2/2022, 11.55/2025} {
  \draw[accentgray] (\x,0.08) -- (\x,-0.08);
  \node[below,font=\scriptsize] at (\x,-0.12) {\yr};
}
\draw[accentgray] (6.4,0.08) -- (6.4,-0.08);
\node[above,font=\scriptsize] at (6.4,0.12) {2024};

\node[yr,minimum width=2.2cm] (cp) at (0.6,1.3) {GitHub\\Copilot};
\node[yr] (cw) at (3.55,1.3) {Amazon\\CodeWhisperer};
\draw[-{Latex[length=2mm]},issatblue] (cp) -- (0.6,0.1);
\draw[-{Latex[length=2mm]},issatblue] (cw) -- (3.2,0.1);
\node[align=center,font=\scriptsize\itshape,text width=3.6cm] at (1.9,2.25) {reactive, line-level\\autocomplete};

\node[ya] (dv) at (6.4,-1.5) {Devin};
\node[ya] (cc) at (9.75,-1.5) {Claude Code\\/ Cursor};
\node[ya] (ox) at (13.1,-1.5) {OpenAI\\Codex CLI};
\draw[-{Latex[length=2mm]},issatgold!80!black] (dv) -- (6.4,-0.1);
\draw[-{Latex[length=2mm]},issatgold!80!black] (cc) -- (9.75,-0.1);
\draw[-{Latex[length=2mm]},issatgold!80!black] (ox) -- (13.1,-0.1);
\node[align=center,font=\scriptsize\itshape,text width=7.2cm] at (9.75,-2.45)
  {autonomous, task-level agents: plan, edit, run tests, revise};
\end{tikzpicture}

\caption{AI coding tools from 2021 to 2025.}
\label{fig:timeline}
\end{figure}

A developer can now hand an agent a feature description and receive a pull request. The informal
term \emph{vibe coding} describes this way of working: the developer states the intent and the agent
writes the implementation. \citet{konrad2026vibe} describe a related effect, \emph{vibe
architecting}, in which structural decisions follow from how a prompt is worded and leave no
recorded rationale.

These agents already contribute to real projects. The AIDev dataset \citep{li2025aidev} collects
pull requests opened by autonomous coding agents on public GitHub repositories, and studies built
on it work with thousands of such requests. Because agent-written code is reviewed and merged into
real software, its structural quality can be studied on real data. This thesis uses the same
dataset.

Uptake and merge figures describe use, which is only one part of software quality. Software
engineering has a longer vocabulary for design: Brooks's analysis of accidental complexity
\citep{brooks1975mythical}, Martin's design principles \citep{martin2000solid}, later grouped under
the acronym SOLID, and the code smells catalogued by \citet{fowler1999refactoring}. Code can pass
its tests and still be tightly coupled and expensive to change, and most of the cost of software is
spent on maintenance after the first release.

Recent studies suggest that AI assistance affects this side of quality. \citet{xu2025productivity}
report a shift in how core developers split their time between writing and reviewing code after
Copilot was introduced, and \citet{zhu2026smells} describe a \emph{Volume-Quality Inverse Law}: in
AI-generated repositories, code volume predicts structural degradation, and neither functional
correctness nor detailed prompting prevents it.

\section{Problem Statement}
\label{sec:problem-statement}

Coding agents are mostly evaluated on properties that are easy to test automatically, and less on
the properties that decide how expensive code will be to maintain.

\paragraph{Benchmarks measure function, not design.} HumanEval \citep{chen2021codex}, MBPP,
BigCodeBench \citep{zhuo2024bigcodebench} and SWE-bench \citep{jimenez2024swebench} all ask whether
generated code passes a set of tests, usually for one function or one issue. A pass rate is well
defined and reproducible, but it says nothing about whether the surrounding code is well structured
or whether a maintainer would accept the change. The distance between the two has been measured:
\citet{metr2026sweverified} found that roughly half of the agent patches that pass SWE-bench Verified
would not be merged by maintainers.

\paragraph{Structural quality is rarely compared across agents.} Studies that do look at the
structure of AI-generated code ask language models to detect SOLID violations
\citep{pehlivan2025solid}, analyse large projects generated by one tool
\citep{beyond2026designissues}, or describe recurring defects in agent-written code
\citep{zhu2026smells}. None of them compares several agents on the same corpus of real
contributions.

\paragraph{The agent matters, not only the model.} Most evaluations vary the language model and
keep the agent software fixed. \citet{bensghaier2026harness} show that the agent software has an
effect of its own, so someone choosing a tool chooses an agent, not only a model. This thesis
compares three agents: Claude Code, GitHub Copilot and OpenAI Codex.

In short, current evaluation practice says little about the design quality of the code agents
contribute to real repositories. Adoption decisions may therefore rest on functional scores alone,
while maintenance costs accumulate elsewhere. The problem addressed by this thesis is whether the
code that coding agents contribute to real projects follows established object-oriented design
principles, in addition to working, and whether the agents differ in this respect.

\section{Research Questions}
\label{sec:research-questions}

The study answers two research questions. Chapter~\ref{ch:methodology}
(Section~\ref{sec:rqs}) describes how each is answered.

\begin{description}[leftmargin=2.4cm,style=nextline]
  \item[RQ1 (Adherence and differentiation)] Do coding agents differ in how well their pull requests
  adhere to the five SOLID principles, and how many of the design problems in the code they touch do
  they introduce themselves, in Java and in Python?

  \item[RQ2 (Measurability)] Can SOLID adherence be measured at repository scale with deterministic
  structural detectors, and how well do these detectors agree with an independent reading judge and
  with established analysers?
\end{description}

\section{Research Objectives}
\label{sec:research-objectives}

The main objective is to find out whether code written by coding agents follows established
object-oriented design principles, in addition to working. The thesis does so in two languages, Java
and Python.

One scoping decision bounds every later claim. The five SOLID principles describe classes and the
relations between them, so all five measurements work below the level of the package. When the
thesis uses the word architecture, it means design structure at this level, not the decomposition of
a system into components and layers.

The objectives are:

\begin{enumerate}[label=\textbf{O\arabic*.}, leftmargin=1.6cm]
  \item Build a reproducible framework that measures adherence to the five SOLID principles (Single
  Responsibility, Open/Closed, Liskov Substitution, Interface Segregation and Dependency Inversion)
  on real pull requests written by coding agents, using one deterministic structural detector per
  principle, in Java and in Python.
  \item Build the corpus from all three outcomes a pull request can have: merged, rejected and still
  open. A corpus of merged requests shows what maintainers accepted, which is not the same as what
  agents produced.
  \item Measure how well each detector agrees with an independent reading of the code. A judge reads
  the source of a stratified sample of flagged and unflagged entities, and precision and recall are
  reported for each principle. The judge is a large language model, a technique validated for SOLID
  detection by \citet{pehlivan2025solid}, so every figure is agreement between the detector and the
  judge.
  \item Compare three agents, Claude Code, GitHub Copilot and OpenAI Codex, with rates reported both
  pooled and averaged over repositories.
  \item Determine where design problems are located: inside one unit (a class or a method) that a
  single change can write in full, or in a relation between units (a dependency, an interface, an
  override) that a single change rarely writes on both sides.
  \item Carry out the whole study in two languages, Java and Python, with the same agents, curation
  rules, validation procedure and statistical model, to find out which results hold in both.
\end{enumerate}

Three design choices frame the study. First, generating many new repositories with each agent would
require compute and API budgets that were not available. The study therefore measures code the
agents actually submitted to open-source projects, taken from the AIDev dataset
\citep{li2025aidev}. This is cheaper, it reflects real tasks, and each pull request is a bounded
change made by one agent.

Second, the study is carried out in two languages, and they play different roles. A common
objection to measuring SOLID is that the principles fit poorly in languages that do not require
object orientation. Java answers that objection: classes are compulsory, \texttt{interface} is a
keyword, \texttt{@Override} is explicit and subtyping is nominal, so each principle has a direct
counterpart in the syntax. Python is the opposite case: classes are optional, contracts are
conventions, and no compiler checks an override against its base. A result found in both languages
is unlikely to be a side effect of Java's type system. Only the fact-extraction layer of the
framework depends on the language, so the Python study uses its own extractor and recalibrated
thresholds and changes nothing else.

Third, the agent set follows from the data. The same three agents are used in both languages so
that every cross-language comparison has the same agents on both sides.
Section~\ref{sec:agent-selection} explains why the other agents were left out.
